\chapter{Performance Evaluation}

\section{Introduction}

This chapter presents the experimental performance evaluation of the implemented system. 
While Chapter 5 provided theoretical timing analysis, this chapter validates system behavior through empirical observation and measurement.

The evaluation focuses on:

\begin{itemize}
	\item Control loop frequency
	\item Communication latency
	\item Actuator responsiveness
	\item Wireless stability
	\item Resource utilization
\end{itemize}

% ------------------------------------------------

\section{Experimental Setup}

The system was tested under standard operating conditions:

\begin{itemize}
	\item Bluetooth connection established with Android control application
	\item Motors connected under nominal load
	\item Servo connected for steering control
	\item Telemetry transmission enabled
\end{itemize}

Measurements were performed using:

\begin{itemize}
	\item Serial debug timestamps
	\item Oscilloscope verification for PWM signals (if available)
	\item Software-based timing counters
\end{itemize}

% ------------------------------------------------

\section{Control Loop Frequency Measurement}

The control loop frequency was measured by recording timestamps at the beginning of each iteration.

The average loop period is defined as:

\[
T_{avg} = \frac{\sum_{i=1}^{N} (t_i - t_{i-1})}{N}
\]

The loop frequency is then:

\[
f_{loop} = \frac{1}{T_{avg}}
\]

\subsection{Observed Results}

Under normal operation:

\begin{itemize}
	\item Loop frequency remained stable above the minimum requirement for RC responsiveness
	\item No blocking delays were observed
	\item Frequency variations were minimal
\end{itemize}

This confirms deterministic execution behavior.

% ------------------------------------------------

\section{End-to-End Command Latency}

End-to-end latency was measured from:

\begin{itemize}
	\item User input on the mobile device
	\item To observable actuator response
\end{itemize}

Latency components include:

\begin{itemize}
	\item Bluetooth transmission delay
	\item Packet parsing delay
	\item Control loop synchronization delay
\end{itemize}

\subsection{Observed Behavior}

Testing indicates:

\begin{itemize}
	\item No perceptible delay in manual control
	\item Immediate motor and servo reaction
	\item Stable control during rapid command changes
\end{itemize}

Measured latency remained within acceptable bounds for real-time manual control applications.

% ------------------------------------------------

\section{Actuator Response Stability}

PWM signals were verified for:

\begin{itemize}
	\item Consistent duty cycle generation
	\item Stable servo pulse width
	\item Absence of jitter-induced irregularities
\end{itemize}

Hardware-assisted PWM ensured stable actuator control independent of minor loop timing variations.

% ------------------------------------------------

\section{Wireless Communication Stability}

Long-duration tests were performed to evaluate:

\begin{itemize}
	\item Packet loss occurrence
	\item Connection drop frequency
	\item Recovery behavior
\end{itemize}

Results show:

\begin{itemize}
	\item Stable Bluetooth connectivity within expected operating range
	\item Proper handling of packet validation
	\item No firmware crashes during communication stress
\end{itemize}

% ------------------------------------------------

\section{Telemetry Throughput}

Telemetry packet transmission rate was evaluated to ensure it does not interfere with control responsiveness.

Observations include:

\begin{itemize}
	\item Telemetry packets transmitted without saturating communication channel
	\item Control commands remain prioritized
	\item No buffer overflows detected
\end{itemize}

This confirms appropriate bandwidth allocation.

% ------------------------------------------------

\section{Memory Utilization}

Memory usage was analyzed based on:

\begin{itemize}
	\item Static memory allocation
	\item Buffer sizes
	\item Stack usage
\end{itemize}

Because dynamic allocation is avoided:

\begin{itemize}
	\item Memory usage remains predictable
	\item Fragmentation risks are eliminated
\end{itemize}

The firmware operates well within the available memory limits of the microcontroller.

% ------------------------------------------------

\section{CPU Utilization Estimation}

CPU utilization can be approximated as:

\[
U = \frac{T_{loop}}{T_{available}}
\]

Given the high processing capability of the microcontroller relative to the task complexity:

\begin{itemize}
	\item CPU load remains low
	\item Significant computational margin exists
	\item Additional features can be integrated without compromising stability
\end{itemize}

% ------------------------------------------------

\section{Stress Testing}

Stress tests included:

\begin{itemize}
	\item Rapid command changes
	\item Continuous telemetry transmission
	\item Simultaneous motor and servo operation
\end{itemize}

The system maintained:

\begin{itemize}
	\item Stable loop frequency
	\item No watchdog resets
	\item No observable performance degradation
\end{itemize}

% ------------------------------------------------

\section{Performance Limitations}

Identified performance constraints include:

\begin{itemize}
	\item Bluetooth bandwidth limitations
	\item I2C transaction latency
	\item Physical actuator response time
\end{itemize}

These limitations are external to core firmware execution and do not indicate architectural flaws.

% ------------------------------------------------

\section{Summary}

Experimental evaluation confirms:

\begin{itemize}
	\item Deterministic control loop execution
	\item Acceptable end-to-end latency
	\item Stable actuator control
	\item Reliable wireless communication
	\item Low resource utilization
\end{itemize}

The measured performance aligns with the theoretical analysis presented in Chapter 5 and validates the effectiveness of the modular firmware architecture.